% section_V_scale_insert.tex -- replaces the n=5 evaluation paragraphs in Sections V and VI.
% Needs: \usepackage{booktabs}.  The \input files are produced by  python3 gen_paper_tables.py  (every number is computed from raw.csv).
% Set \texpath to where the generated files live, e.g.  \newcommand{\texpath}{results/paper_tex}
\providecommand{\texpath}{results/paper_tex}
\providecommand{\inputif}[1]{\IfFileExists{#1}{\input{#1}}{\fbox{missing generated file #1}}}

\subsection{Sample size, path variance and statistics}
\label{sec:scale}
Earlier tables in this paper used $n=5$ runs per cell on hand-picked paths. That is too few to support the percentile columns
(with $n=5$ the p95 and p99 are both essentially the sample maximum) and says nothing about how results move with the path. We therefore
re-ran the transfer comparison as a multi-seed, multi-path matrix. Three \emph{path families} are defined as parameter ranges
(Table~\ref{tab:families}); seed $s\in\{0,\dots,29\}$ draws one concrete path from each family, and every policy is run on exactly that path
(a paired design) in a seed-dependent random order, so slow drift in machine load cannot favour one policy. Each cell therefore has
$n=30$ independent runs of a 1\,MB transfer, each on a different path, and the interval estimates include both path variance and run noise.

\begin{table}[t]\centering\footnotesize
\caption{Path families (one-way delay; the proxy adds the same delay to ACKs). Each seed draws uniformly from the ranges. Cross traffic is a synthetic on/off process at the shared bottleneck (proxy backend only).}
\label{tab:families}
\begin{tabular}{@{}lp{5.6cm}@{}}\toprule
family & ranges\\\midrule
dc\_microburst & 0.25--1\,ms; 8--12\,MB/s; 20--40\,kB buffer; no random loss; 1--3\,ms line-rate bursts every 20--40\,ms\\
enterprise\_campus & 2--8\,ms, 0--2\,ms jitter; 2.5--5\,MB/s; 60--120\,kB buffer; 0--0.1\% loss; on/off cross traffic\\
cellular\_highlat & 40--120\,ms, 10--40\,ms jitter; 150--600\,kB/s; 20--50\,kB buffer; 0.5--3\% loss\\
\bottomrule\end{tabular}
\end{table}

For each cell we report the mean with a $t$-based 95\% confidence interval, the median and p95 with 95\% percentile-bootstrap intervals
(4000 resamples), and p99 as a point estimate only: with $n=30$ it interpolates between the two largest observations and has no meaningful interval.
Policies are compared by the \emph{ratio} of completion times on the same seed; we report the geometric mean ratio with a 95\% interval from a
$t$-interval on the log ratios, a paired $t$-test and, in the CSV, a Wilcoxon signed-rank test, with Holm correction within each family and reference.
Runs that fail the byte-for-byte check are counted ($n_{ok}/n$) and excluded from the timing statistics; a pair is dropped when either member failed.

\subsection{Results: UDP-impairment-proxy backend ($n=30$ per cell)}
\inputif{\texpath/proxy_scale_table.tex}
\inputif{\texpath/proxy_scale_paired.tex}
\inputif{\texpath/proxy_scale_text.tex}

\subsection{Results: kernel TCP under Linux netem ($n=30$ per cell)}
\label{sec:kerneltcp}
This backend replaces the user-space proxy with \texttt{tc netem} and \texttt{tbf} on the loopback device and adds native kernel TCP (CUBIC and BBR) as
baselines. Kernel TCP sends the same random payload, over the same bottleneck, with its timer started after \texttt{connect()} (the REPLIC timer starts after
its handshake) and stopped when the receiver has all bytes and answers, which is the same end condition as the REPLIC client. Netem impairs both directions of
\texttt{lo}, whereas the proxy impairs the forward path only, so \textbf{the two backends are never mixed in one table}. The synthetic cross-traffic bursts of the proxy have
no netem counterpart; under netem \texttt{dc\_microburst} keeps only its shallow buffer, delay and rate.
\inputif{\texpath/netem_scale_table.tex}
\inputif{\texpath/netem_scale_paired.tex}
\inputif{\texpath/netem_scale_text.tex}
\inputif{\texpath/netem_matrix.tex}

\subsection{Coexistence with $n=30$ runs per cell}
\label{sec:fairness30}
The fairness cells of Section~\ref{sec:fairness} were repeated with 30 runs each (subject against $k$ CUBIC or BBR background flows, $k\in\{1,2\}$).
\begin{table*}[t]\centering\footnotesize
\caption{Coexistence, mean $\pm$ s.d. over 30 runs per row. $J$: Jain index over all flows; ``share'': subject throughput divided by the equal share; ``min'': smallest flow's share divided by the equal share.}
\label{tab:fair30}
\begin{tabular}{@{}llrrrrr@{}}\toprule
subject & background & $k$ & $n$ & $J$ & share & min\\\midrule
\input{\texpath/fairness30_rows.tex}\\
\bottomrule\end{tabular}
\end{table*}

\paragraph{Reading Table~\ref{tab:fair30}.} With 30 runs per cell (Welch tests, unadjusted) the REPLIC agent has a higher Jain index than the static agent against CUBIC
($k=1$: 0.974 vs 0.942, $p=0.012$; $k=2$: 0.958 vs 0.881, $p<0.001$), and than the static agent with the \texttt{--polite} guard ($k=1$: $p=0.001$; $k=2$: 0.958 vs 0.919, $p=0.015$);
at $n=5$ the guard comparison was not significant. Against BBR the picture is different: REPLIC's Jain index is not distinguishable from the static agent's
($k=1$: $p=0.34$; $k=2$: $p=0.33$), and it takes only about $0.74\times$ its equal share, i.e.\ it yields to the simplified BBR background flows rather than being unfair to them,
at the cost of its own throughput. We make no fairness claim against BBR beyond ``not a bully''.

\subsection{Threats to validity}
(i) The path families are our own parameter ranges, not measurements of deployed networks. (ii) Intervals describe sampling error over our families, not any particular real path.
(iii) On \texttt{dc\_microburst} a 1\,MB transfer ends within a few round-trips, so handshake and slow start, which are fixed rules shared by all REPLIC policies, dominate
the comparison. (iv) The learned policies were trained on a simulator with simplified CUBIC and BBR; no result here should be read as evidence about real BBR deployments.
(v) A single machine with two cores runs sender, receiver and bottleneck emulation; the order randomisation limits, but does not remove, load-dependent timing noise.
(vi) Kernel-TCP retransmission counts come from a system-wide counter and are only meaningful on an otherwise idle host.
